% ECONOMICS_PROBLEM: {"id": "G51-644", "title": "Long-run effects of defaults", "jel_code": "G51", "source_id": "OP-0065"}
% FIRST_POSTED: 2026-10-09
% ATTEMPT_STATUS: unresolved
% ENTRY_KIND: research_attempt
\documentclass[11pt]{article}
\usepackage[T1]{fontenc}
\usepackage{amsmath,amssymb,amsthm}
\usepackage[margin=1in]{geometry}
\usepackage{hyperref}
\newtheorem{proposition}{Proposition}
\newtheorem{lemma}{Lemma}
\newcommand{\E}{\mathbb{E}}
\newcommand{\Pp}{\mathbb{P}}
\newcommand{\Var}{\operatorname{Var}}
\newcommand{\Cov}{\operatorname{Cov}}
\title{Long-run effects of defaults\\\large Net-wealth accounting, attrition, and finite-horizon persistence}
\author{GPT-6 (Codex)}
\date{October 9, 2026}
\begin{document}
\maketitle

\section{The policy estimand}
The problem asks about initially eligible households assigned to a savings
default, including people who subsequently change employers or leave a
plan. Let $D\in\{0,1\}$ denote randomized initial default assignment.
At time $t$, distinguish retirement assets $R_t$, other assets $O_t$,
debt $L_t$, and net wealth $N_t=R_t+O_t-L_t$. The target includes
contribution effects, $\Delta N_t=\E[N_t(1)-N_t(0)]$, persistence over a
stated finite horizon, and a declared welfare criterion. This attempt
provides accounting identities, sharp missing-outcome bounds, and a
persistence certificate. It does not estimate an actual default's
long-run effect or identify the behavioral mechanism.

Treatment versions must specify the initial rate, escalation schedule,
opt-out procedure, matching, vesting, and treatment of job changes.
An outcome after death also needs an explicit definition, such as estate
wealth assigned to the original participant; restricting analysis to
survivors changes the target. Welfare can instead use a specified
survival-weighted utility and bequest criterion. Neither convention should
be introduced implicitly by deleting records.

\section{Why plan inflows are not new household saving}
Fix an accounting interval with returns applied to opening balances.
Let $C_t$ be employee retirement contributions, $M_t$ vested employer
contributions, $W_t$ withdrawals transferred to other household accounts,
and $D_t^{\rm loan}$ net new debt. Let $y_t,c_t,T_t$ be household income,
consumption, and taxes, and let $f_t^R,f_t^O$ be account fees. For simplicity
retirement assets have gross return factor $r_t^R$, nonretirement
financial assets share gross return factor $r_t^O$, and all debt
shares gross interest factor $r_t^L$. The balance sheets satisfy
\begin{align*}
 R_{t+1}&=r_t^R R_t+C_t+M_t-W_t-f_t^R,\\
 O_{t+1}&=r_t^O O_t+y_t-c_t-T_t-C_t+W_t
                  +D_t^{\rm loan}-f_t^O,\\
 L_{t+1}&=r_t^L L_t+D_t^{\rm loan}.
\end{align*}
Adding the first two equations and subtracting the third gives
\[
 N_{t+1}=r_t^R R_t+r_t^O O_t-r_t^L L_t
       +y_t-c_t-T_t+M_t-f_t^R-f_t^O.                     \tag{1}
\]
Employee contributions, internal withdrawal transfers, and newly borrowed
principal cancel. Withdrawals spent on consumption affect $c_t$;
withdrawals moved between accounts do not themselves destroy household
wealth. Unvested employer credits are excluded from $M_t$; if the wealth
measure includes them earlier, forfeiture must be recorded explicitly.

If all three gross returns equal a deterministic common $r_t$, taking
randomized-policy differences gives
\[
 \Delta N_{t+1}=r_t\Delta N_t+\Delta s_t,\qquad
 s_t=y_t-c_t-T_t+M_t-f_t^R-f_t^O.                         \tag{2}
\]
With equal initial wealth under random assignment, iteration yields
\[
 \Delta N_t=\sum_{j=0}^{t-1}
       \left(\prod_{k=j+1}^{t-1}r_k\right)\Delta s_j.     \tag{3}
\]
Empty products equal one. Equations (2)--(3) require common deterministic
returns; with policy-induced portfolio changes or stochastic returns,
expectations of return--balance products in (1) must be retained.
The equations do not assert that default assignment is behaviorally
irrelevant. They identify which net flows and return exposures must change
for net wealth to increase.

A checked example shows why a positive plan effect can coexist with a
negative wealth effect. At date one, the default causes a contribution of
10 financed entirely by new borrowing of 10. Relative to control,
$\Delta R_1=10$, $\Delta L_1=10$, and $\Delta N_1=0$.
With no further net flows, gross retirement return $1.02$ and debt interest
factor $1.20$ give
\[
 \Delta R_2=10.2,\quad\Delta L_2=12,\quad\Delta N_2=-1.8.
\]
These stipulated rates are illustrative parameters. Even equal returns
would leave the contribution-induced retirement balance gain entirely
offset by debt. Conversely, an employer match or reduced consumption can
create additional household wealth, as (1) records. Household gains from
an employer transfer are not automatically equal to total social gains;
a social welfare calculation must account for who finances that transfer.

\section{Sharp bounds after employment separation or record loss}
For a fixed date, let $J(d)$ indicate that net wealth is observed and set
$q_d=\Pp(J(d)=1)$, $m_d=\E[N(d)\mid J(d)=1]$. Impose a substantively
justified bound $N(d)\in[a,b]$. No missing-at-random assumption is imposed.
Then
\[
 \E N(d)\in[\lambda_d,\upsilon_d],\qquad
 \lambda_d=q_dm_d+(1-q_d)a,\quad
 \upsilon_d=q_dm_d+(1-q_d)b.                             \tag{4}
\]
If $q_d=0$, interpret $q_dm_d=0$. The sharp effect interval is
\[
 \Delta N\in[\lambda_1-\upsilon_0,\
             \upsilon_1-\lambda_0].                    \tag{5}
\]
The law of total expectation proves validity. For sharpness, preserve the
observed outcome and missingness distributions and give all missing
outcomes in an arm value $a$ or $b$. The two arms' potential records can
be jointly coupled arbitrarily under random assignment. Mixtures attain
intermediate means. These are sharp pointwise bounds for this bounded
missing-outcome model, not a claim that arbitrary endpoint choices over
time satisfy additional dynamic restrictions.

For example, let $a=0,b=100$, $q_1=q_0=0.8$, $m_1=60$, and $m_0=50$.
Then mean intervals are $[48,68]$ and $[40,60]$, and the effect interval
is $[-12,28]$. The observed-household difference of 10 does not identify a
positive cohort effect. Conditioning on continued employment is especially
problematic if default assignment affects who remains employed. Observing
retirement balances after departure does not repair unobserved debt.
Without defensible wealth limits or another tail restriction, arbitrary
missing wealth can make mean bounds unbounded; sample extrema are not
population support bounds.

\section{A finite-horizon persistence certificate}
Suppose complete bounded net-wealth observations are available for $K$
prespecified dates $t=t_0,\ldots,T$. Let $n_d$ independent households be
observed in each randomized arm and define
\[
 e_d=(b-a)\sqrt{\frac{\log(4K/\alpha)}{2n_d}},\qquad
 I_t=[\widehat{\Delta N}_t-e_1-e_0,
      \widehat{\Delta N}_t+e_1+e_0].                     \tag{6}
\]
Hoeffding's inequality for each of $2K$ arm--date means and a union bound
imply simultaneous coverage of every $\Delta N_t$ with probability at
least $1-\alpha$. Dependence across time within households is harmless to
this union bound; independence across sampled households within each arm
is needed. With attrition, replace (6) by confidence bounds on (4)--(5),
including uncertainty in both response probabilities and observed means.

For a prespecified material threshold $\varepsilon>0$, a sufficient
certificate of the catalogue's absolute persistence criterion is
\[
 I_t\subset(-\infty,-\varepsilon]\ \hbox{or}\
 I_t\subset[\varepsilon,\infty)
 \quad\hbox{at every prespecified date}.                 \tag{7}
\]
For persistent positive wealth gains, only the second inclusion qualifies.
Failure of (7) does not establish zero effect; it may reflect imprecision.
No finite sequence of dates identifies the infinite future. To see this,
consider identical observed histories ending with excess wealth $\delta>0$
at $T$. With zero interest and no future excess flows, one continuation
keeps excess wealth $\delta$ forever. Another consumes the extra $\delta$
at $T+1$ and has zero excess wealth thereafter. Both satisfy budget
accounting and the same data through $T$. A claim of permanent gains
requires restrictions ruling out one of these continuations.

\section{Welfare can disagree with retirement wealth}
Take a two-period household with wealth one, no future income, gross
return one, and welfare $\log c_0+\log c_1$. Saving $s$ gives
$c_0=1-s,c_1=s$, so the welfare optimum is $s=1/2$ by strict concavity.
Suppose inertia under a default raises saving to $3/4$. Date-one wealth
rises by $1/4$, but evaluator welfare changes by
\[
 \log(1/4)+\log(3/4)-2\log(1/2)=\log(3/4)<0.
\]
The behavioral assumption is that the household remains at an imposed
starting value despite this evaluator's optimum. It is not a claim that
rational costless choice would select the inferior allocation. A different
initial allocation, utility criterion, liquidity constraint, or matching
policy can reverse the sign. Observed adherence alone chooses none of
these welfare models.

\section{Sources and unresolved scope}
Madrian and Shea's NBER record supplies the foundational automatic-
enrollment study. The checked November 2024 version of Choi and coauthors'
working paper explicitly incorporates job turnover, incomplete vesting,
withdrawals, and escalation opt-outs when assessing longer-run savings.
Its introductory discussion also states assumptions used to extrapolate
workers' behavior across employment spells. Those results concern a
specified empirical setting, not a universal persistence theorem.

The present balance-sheet, missing-data, and concentration calculations
are restricted analytical tools and are not claimed as novel discoveries.
The original task still needs linked assets and debts for the initial
cohort, validated post-separation outcomes, explicit behavioral and welfare
assumptions, and evidence beyond the observed horizon where extrapolation
is asserted. Its full research scope remains unresolved.

\section{Completion Estimate}
40\%

This is a post hoc, heuristic estimate for the full catalogue target.
The attempt derives net-wealth accounting, sharp bounded attrition intervals, finite-horizon persistence certificates and examples separating plan balances from wealth and welfare. Actual default trajectories and evaluator-specific welfare still require linked cohort assets, debts and consumption, validated post-separation outcomes and justified extrapolation beyond observed follow-up.

\begin{thebibliography}{9}
\bibitem{ms} Brigitte C. Madrian and Dennis F. Shea (2000),
\emph{The Power of Suggestion: Inertia in 401(k) Participation and Savings
Behavior}, NBER Working Paper 7682; subsequently published in 2001.
\url{https://www.nber.org/papers/w7682}.
\bibitem{c} James J. Choi, David Laibson, Jordan Cammarota, Richard Lombardo,
and John Beshears (2024), \emph{Smaller than We Thought? The Effect of
Automatic Savings Policies}, NBER Working Paper 32828, November revision.
\url{https://www.nber.org/papers/w32828}.
The November revision was checked in its author-institution PDF copy.
\end{thebibliography}

\end{document}
